% CPU pipeline timing: single-issue, superscalar and multithreaded (Lesson 1).
% Three charts of a five-stage pipeline (F D S E W).  (a) one instruction
% enters per cycle; (b) a two-wide superscalar enters two per cycle;
% (c) the same two-wide pipeline fed by four threads (one shade per thread).
% Redraws the outline figures m1l2_cpupipeline and m1l2_multi_thread.
% Usage: % CPU pipeline timing: single-issue, superscalar and multithreaded (Lesson 1).
% Three charts of a five-stage pipeline (F D S E W).  (a) one instruction
% enters per cycle; (b) a two-wide superscalar enters two per cycle;
% (c) the same two-wide pipeline fed by four threads (one shade per thread).
% Redraws the outline figures m1l2_cpupipeline and m1l2_multi_thread.
% Usage: % CPU pipeline timing: single-issue, superscalar and multithreaded (Lesson 1).
% Three charts of a five-stage pipeline (F D S E W).  (a) one instruction
% enters per cycle; (b) a two-wide superscalar enters two per cycle;
% (c) the same two-wide pipeline fed by four threads (one shade per thread).
% Redraws the outline figures m1l2_cpupipeline and m1l2_multi_thread.
% Usage: % CPU pipeline timing: single-issue, superscalar and multithreaded (Lesson 1).
% Three charts of a five-stage pipeline (F D S E W).  (a) one instruction
% enters per cycle; (b) a two-wide superscalar enters two per cycle;
% (c) the same two-wide pipeline fed by four threads (one shade per thread).
% Redraws the outline figures m1l2_cpupipeline and m1l2_multi_thread.
% Usage: \input{figures/l01-cpu-pipeline}
% Width: body (120 mm); the drawing is about 118 mm wide.
\begin{tikzpicture}[x=3.2mm, y=3.2mm, font=\scriptsize\sffamily]
  % \pipe{x}{y}{fill}: one instruction, five cells from (x,y)
  \newcommand{\pipe}[3]{%
    \foreach \s [count=\i from 0] in {F,D,S,E,W} {
      \draw[fill=#3] (#1+\i,#2) rectangle ++(1,-1)
        node[midway] {\s};
    }}
  % (a) single issue
  \node[anchor=south west, font=\scriptsize\sffamily\bfseries] at (0,0.3)
    {\iflangja{(a) 単一発行}{(a) single issue}};
  \foreach \k in {0,...,3} { \pipe{\k}{-\k}{ln-box} }
  \draw[->, ln-muted] (0,-4.6) -- (8,-4.6)
    node[below left, inner sep=1pt] {\iflangja{サイクル}{cycle}};
  % (b) superscalar, two wide
  \node[anchor=south west, font=\scriptsize\sffamily\bfseries] at (13,0.3)
    {\iflangja{(b) 2命令幅スーパースカラ}{(b) two-wide superscalar}};
  \foreach \k in {0,...,3} {
    \pipe{13+\k}{-2*\k}{ln-box}
    \pipe{13+\k}{-2*\k-1}{ln-box}
  }
  \draw[->, ln-muted] (13,-8.6) -- (21,-8.6)
    node[below left, inner sep=1pt] {\iflangja{サイクル}{cycle}};
  % (c) multithreading on the same two-wide pipeline
  \node[anchor=south west, font=\scriptsize\sffamily\bfseries] at (26,0.3)
    {\iflangja{(c) マルチスレッディング}{(c) multithreading}};
  \foreach \k/\c in {0/ln-accent!25, 1/ln-tip!25, 2/ln-caution!20, 3/ln-source!20} {
    \pipe{26+\k}{-2*\k}{\c}
    \pipe{26+\k}{-2*\k-1}{\c}
  }
  \draw[->, ln-muted] (26,-8.6) -- (34,-8.6)
    node[below left, inner sep=1pt] {\iflangja{サイクル}{cycle}};
  \node[anchor=north west, text=ln-muted, align=left] at (26,-9.3)
    {\iflangja{色＝スレッド}{one shade = one thread}};
\end{tikzpicture}

% Width: body (120 mm); the drawing is about 118 mm wide.
\begin{tikzpicture}[x=3.2mm, y=3.2mm, font=\scriptsize\sffamily]
  % \pipe{x}{y}{fill}: one instruction, five cells from (x,y)
  \newcommand{\pipe}[3]{%
    \foreach \s [count=\i from 0] in {F,D,S,E,W} {
      \draw[fill=#3] (#1+\i,#2) rectangle ++(1,-1)
        node[midway] {\s};
    }}
  % (a) single issue
  \node[anchor=south west, font=\scriptsize\sffamily\bfseries] at (0,0.3)
    {\iflangja{(a) 単一発行}{(a) single issue}};
  \foreach \k in {0,...,3} { \pipe{\k}{-\k}{ln-box} }
  \draw[->, ln-muted] (0,-4.6) -- (8,-4.6)
    node[below left, inner sep=1pt] {\iflangja{サイクル}{cycle}};
  % (b) superscalar, two wide
  \node[anchor=south west, font=\scriptsize\sffamily\bfseries] at (13,0.3)
    {\iflangja{(b) 2命令幅スーパースカラ}{(b) two-wide superscalar}};
  \foreach \k in {0,...,3} {
    \pipe{13+\k}{-2*\k}{ln-box}
    \pipe{13+\k}{-2*\k-1}{ln-box}
  }
  \draw[->, ln-muted] (13,-8.6) -- (21,-8.6)
    node[below left, inner sep=1pt] {\iflangja{サイクル}{cycle}};
  % (c) multithreading on the same two-wide pipeline
  \node[anchor=south west, font=\scriptsize\sffamily\bfseries] at (26,0.3)
    {\iflangja{(c) マルチスレッディング}{(c) multithreading}};
  \foreach \k/\c in {0/ln-accent!25, 1/ln-tip!25, 2/ln-caution!20, 3/ln-source!20} {
    \pipe{26+\k}{-2*\k}{\c}
    \pipe{26+\k}{-2*\k-1}{\c}
  }
  \draw[->, ln-muted] (26,-8.6) -- (34,-8.6)
    node[below left, inner sep=1pt] {\iflangja{サイクル}{cycle}};
  \node[anchor=north west, text=ln-muted, align=left] at (26,-9.3)
    {\iflangja{色＝スレッド}{one shade = one thread}};
\end{tikzpicture}

% Width: body (120 mm); the drawing is about 118 mm wide.
\begin{tikzpicture}[x=3.2mm, y=3.2mm, font=\scriptsize\sffamily]
  % \pipe{x}{y}{fill}: one instruction, five cells from (x,y)
  \newcommand{\pipe}[3]{%
    \foreach \s [count=\i from 0] in {F,D,S,E,W} {
      \draw[fill=#3] (#1+\i,#2) rectangle ++(1,-1)
        node[midway] {\s};
    }}
  % (a) single issue
  \node[anchor=south west, font=\scriptsize\sffamily\bfseries] at (0,0.3)
    {\iflangja{(a) 単一発行}{(a) single issue}};
  \foreach \k in {0,...,3} { \pipe{\k}{-\k}{ln-box} }
  \draw[->, ln-muted] (0,-4.6) -- (8,-4.6)
    node[below left, inner sep=1pt] {\iflangja{サイクル}{cycle}};
  % (b) superscalar, two wide
  \node[anchor=south west, font=\scriptsize\sffamily\bfseries] at (13,0.3)
    {\iflangja{(b) 2命令幅スーパースカラ}{(b) two-wide superscalar}};
  \foreach \k in {0,...,3} {
    \pipe{13+\k}{-2*\k}{ln-box}
    \pipe{13+\k}{-2*\k-1}{ln-box}
  }
  \draw[->, ln-muted] (13,-8.6) -- (21,-8.6)
    node[below left, inner sep=1pt] {\iflangja{サイクル}{cycle}};
  % (c) multithreading on the same two-wide pipeline
  \node[anchor=south west, font=\scriptsize\sffamily\bfseries] at (26,0.3)
    {\iflangja{(c) マルチスレッディング}{(c) multithreading}};
  \foreach \k/\c in {0/ln-accent!25, 1/ln-tip!25, 2/ln-caution!20, 3/ln-source!20} {
    \pipe{26+\k}{-2*\k}{\c}
    \pipe{26+\k}{-2*\k-1}{\c}
  }
  \draw[->, ln-muted] (26,-8.6) -- (34,-8.6)
    node[below left, inner sep=1pt] {\iflangja{サイクル}{cycle}};
  \node[anchor=north west, text=ln-muted, align=left] at (26,-9.3)
    {\iflangja{色＝スレッド}{one shade = one thread}};
\end{tikzpicture}

% Width: body (120 mm); the drawing is about 118 mm wide.
\begin{tikzpicture}[x=3.2mm, y=3.2mm, font=\scriptsize\sffamily]
  % \pipe{x}{y}{fill}: one instruction, five cells from (x,y)
  \newcommand{\pipe}[3]{%
    \foreach \s [count=\i from 0] in {F,D,S,E,W} {
      \draw[fill=#3] (#1+\i,#2) rectangle ++(1,-1)
        node[midway] {\s};
    }}
  % (a) single issue
  \node[anchor=south west, font=\scriptsize\sffamily\bfseries] at (0,0.3)
    {\iflangja{(a) 単一発行}{(a) single issue}};
  \foreach \k in {0,...,3} { \pipe{\k}{-\k}{ln-box} }
  \draw[->, ln-muted] (0,-4.6) -- (8,-4.6)
    node[below left, inner sep=1pt] {\iflangja{サイクル}{cycle}};
  % (b) superscalar, two wide
  \node[anchor=south west, font=\scriptsize\sffamily\bfseries] at (13,0.3)
    {\iflangja{(b) 2命令幅スーパースカラ}{(b) two-wide superscalar}};
  \foreach \k in {0,...,3} {
    \pipe{13+\k}{-2*\k}{ln-box}
    \pipe{13+\k}{-2*\k-1}{ln-box}
  }
  \draw[->, ln-muted] (13,-8.6) -- (21,-8.6)
    node[below left, inner sep=1pt] {\iflangja{サイクル}{cycle}};
  % (c) multithreading on the same two-wide pipeline
  \node[anchor=south west, font=\scriptsize\sffamily\bfseries] at (26,0.3)
    {\iflangja{(c) マルチスレッディング}{(c) multithreading}};
  \foreach \k/\c in {0/ln-accent!25, 1/ln-tip!25, 2/ln-caution!20, 3/ln-source!20} {
    \pipe{26+\k}{-2*\k}{\c}
    \pipe{26+\k}{-2*\k-1}{\c}
  }
  \draw[->, ln-muted] (26,-8.6) -- (34,-8.6)
    node[below left, inner sep=1pt] {\iflangja{サイクル}{cycle}};
  \node[anchor=north west, text=ln-muted, align=left] at (26,-9.3)
    {\iflangja{色＝スレッド}{one shade = one thread}};
\end{tikzpicture}
